%% THERMODYNAMIQUE %%
%
\section{Temp�rature pseudo-adiabatique du thermom�tre mouill�}

\itemstyle{Op�rateur~:} {\sf TPrimew}
\index{TPrimew}

%%%%%OBJET%%%%
\itemstyle{Objet~:}
Temp�rature pseudo-adiabatique du thermometre mouill�
d'une particule donn�e d�finie par : pression, temp�rature, humidit�.

%%%%%ARGUMENTS%%%%
\bigskip
%%%%%ENTREES%%%%
\begin{tabarg}
\multicolumn{3}{l}{\titretab {Entr\'ees}}\\
\hline
$P_s$ & Vecteur, {\tt float} & Pression de la particule d'air (hPa)\\
$T_s$ & Vecteur, {\tt float} & Temp�rature de la particule d'air ($K$ ou $^\circ C$)\\
$Hum$ & Vecteur, {\tt float} & Humidit� relative de la particule d'air (\%)\\
\end{tabarg}

\medskip
%%%%%SORTIES%%%%
\begin{tabarg}
\multicolumn{3}{l}{\titretab {Sorties}}\\
\hline
$T'_w$ & Vecteur, {\tt float} & Temp�rature pseudo-adiabatique du thermom�tre mouill� (m\^eme unit� que $T_s$ )\\
\end{tabarg}


%%%%%FORMULE%%%%
\itemstyle{Formule}

%%TPrimew
% Calcul de la temp�rature pseudo-adiabatique du thermom�tre mouill�
% � partir de Ps, Ts et Hum 

La temp�rature pseudo-adiabatique du thermom�tre mouill� est d�duite
du point de condensation.
Le point de condensation ($P_{condensation}$ et $T_{condensation}$) 
est calcul� par les fonctions $pcond()$ et $tcond()$ .

Puis, la temp�rature pseudo-adiabatique au niveau de pression $P_s$
( ou temp�rature pseudo-adiabatique du thermom�tre mouill� ou plus
commun�ment point de courbe bleue ) est calcul�e par la fonction $twcurv()$ .

Si la temp�rature $T_s$ est inf�rieure � -40 degr�s Celsius, la temp�rature
pseudo-adiabatique du thermom�tre mouill� est alors mise � donn�e manquante.



\itemstyle{Commentaire}

Les fonctions $pcond()$ , $tcond()$ et $twcurv()$ sont d�crites en annexe
\ref{annexe-thermo}.

%%%%%AUTEUR%%%%
\itemstyle{Auteur~:}
CNRM/GMEI/TRAMM (d'apr�s biblioth�que LIBMTO D.Puech)

%%%%%REFERENCE%%%%
\itemstyle{R\'ef\'erences~:}
Triplet-Roche.
